% ===================================================================
%  TAFEL Nr. 6 — Das Haus von innen. --- Möbel. --- Essen. --- Licht.
%  Traduction 2026 d'apres texto/io, controlee sur texto/fr et
%  texto/en. Consigne : voir texto/de/10-tafel-01.tex.
%
%  Deux notes, toutes deux marquees « (*) », et deux appels : celui de
%  « babo » au bloc c1-12-1, celui de « lambrequino » au bloc c2-01-1.
%  L'ordre les departage, comme dans les autres colonnes.
%
%  LA FEMME DE CHAMBRE N'A PAS DE GROS PLAN. L'ido du bloc c1-06-1
%  porte « la chambristino (6) », le francais « la femme de chambre
%  (16) » : c'est l'ido qui s'est trompe, puisque son propre numero 6
%  est l'eponge du bloc c1-02-1. L'allemand suit donc le francais et
%  ecrit (16), et le renvoi reste du texte ordinaire, faute de savoir
%  ou pointer.
%
%  LE TITRE MEME DE LA PAGE PORTE LA REPONSE DE CE TABLEAU, ET ELLE
%  RENVERSE CE QU'ON ATTEND.
%
%  « MÖBEL » EST UN MOT FRANCAIS — le meuble, le bien MEUBLE du droit,
%  ce qui se deplace par opposition a l'immeuble. Il est entre en
%  allemand au XVIIe siecle. Or PAS UN SEUL des meubles que la page
%  enumere n'est emprunte :
%     Schrank, Bett, Tisch, Stuhl, Bank, Schemel, Truhe, Spiegel,
%     Uhr, Kasten, Bord, Herd, Ofen, Bettgestell, Nachttisch.
%  On attendrait le contraire : que la categorie soit vieille et
%  chaque objet importe. C'est l'inverse. L'ABSTRACTION EST EMPRUNTEE
%  QUAND CHACUN DE SES MEMBRES EST INDIGENE — parce qu'un lit et une
%  table etaient la depuis toujours, tandis que « l'ensemble des biens
%  meubles » est une notion de notaire, arrivee avec le droit et avec
%  l'ebeniste francais.
%
%  C'est le tableau 1 confirme par l'autre bout. Il avait trouve
%  l'abstrait latin et le maniable germanique ; ici l'abstrait ne
%  regroupe QUE du maniable, et il est francais quand meme.
%
%  ET LA NOTE (*) DU LAMBREQUIN DONNE LA PREUVE PAR LE DETAIL.
%  Guignon enumere les langues qui disent toutes « lambrequin » et
%  compte l'allemand parmi elles. C'est vrai, mais ce n'est pas le mot
%  qu'on emploie : le pente de velours de l'alinea 1 s'appelle
%  d'ordinaire Schabracke — et « Schabracke » est le turc
%  « çaprak », arrive par le hongrois avec la housse de selle des
%  cavaliers, pendant les guerres contre les Ottomans. LE PENTE DE
%  VELOURS DU SALON BOURGEOIS PORTE LE NOM D'UNE COUVERTURE DE CHEVAL
%  PRISE A L'ENNEMI. La colonne ecrit « Schabracke » a l'alinea 1,
%  parce que c'est le mot allemand ; la note, elle, garde
%  « Lambrequin », parce que c'est le mot que la note cite.
%
%  ET LA SECONDE NOTE, CELLE DU « babo », OUBLIE L'ALLEMAND. Guignon
%  aligne « baby », « bébé », « bambino » et n'ecrit pas le mot
%  allemand — il n'y en a pas. L'allemand a « Säugling », qui est le
%  nourrisson du medecin, et « Kleinkind », qui est une classe d'age ;
%  pour le bebe qu'on berce, il a emprunte l'anglais « Baby » au XIXe
%  siecle, et c'est le mot de l'alinea 12. Le seul mot de tendresse de
%  la page est un emprunt.
%
%  DEUX INVERSIONS, ET C'EST ENCORE LA MEME. Le champ median de la
%  phrase allemande range ses complements par leur poids et non par
%  l'ordre de l'ido : « den Fingerhut … in den Nähtisch räumen »
%  rejetait la table a ouvrage (38) derriere les quatre objets
%  qu'elle recoit, et « mit einem Gummischlauch an den Gasbrenner »
%  mettait le tuyau (61) avant le bec (62). Le remede est le meme
%  qu'aux tableaux 4 et 5, pris par l'autre bout : porter le
%  complement EN TETE du champ median quand l'ido le met en tete.
%
%  ENFIN LA CUISINE, ALINEAS 1 A 14 DE LA QUATRIEME COLONNE, EST LA
%  PIECE OU L'ALLEMAND A LE PLUS DE MOTS ET LE MOINS D'EMPRUNTS :
%  Topf, Kessel, Pfanne, Schmortopf, Bratpfanne, Tiegel, Zuber, Krug,
%  Kanne, Schale, Napf. Une seule exception, et c'est la plus visible :
%  la Kasserolle des alineas 6 et 14, qui est francaise. Elle
%  est arrivee avec la cuisine bourgeoise, comme le dessert du
%  tableau 4.
% ===================================================================

%%K t06-apar-1 apar
\VUsaut{6.0mm}
\VUtitre{15.0pt}{60}{TAFEL Nr. 6}
\VUsaut{1.4mm}
\VUtitre{11.4pt}{0}{Das Haus von innen, Möbel, Essen,}
\VUsaut{0.4mm}
\VUtitre{11.4pt}{0}{Licht.}
\VUsaut{2.2mm}

%%K t06-apar-2 apar
Das Haus ist fertig. Johanns Onkel ist in sein neues Heim gezogen, dessen
Einteilung wir schon kennen. Sehen wir uns an, wie er es eingerichtet hat. Zuerst
besuchen wir:

%%K t06-tit-1 sub
\VUpk{10.2pt}{40}{Das Schlafzimmer.}
\VUsaut{3.0mm}

%%K t06-c1-01-1 p
1. --- Wir kommen ungelegen: das Mädchen macht gerade das \VUgras{Zimmer} sauber.

%%K t06-c1-02-1 p
2. --- Auf dem \VUgras{Waschtisch} \textsuperscript{(1)} liegen verstreut die Dinge,
mit denen man sich wäscht, kämmt und fertig macht. Es sind die
\VUgras{Waschschüssel} \textsuperscript{(2)} und der \VUgras{Wasserkrug}
\textsuperscript{(3)}, die \VUgras{Haarbürste} \textsuperscript{(4)}, der
\VUgras{Kamm} \textsuperscript{(5)}, die \VUgras{Seife} \textsuperscript{(6)} und der
\VUgras{Schwamm} \textsuperscript{(7)}, und zuletzt ein \VUgras{Fläschchen}
\textsuperscript{(8)} Mundwasser.

%%K t06-c1-03-1 p
3. --- Auf dem Boden, vor dem Waschtisch, steht der \VUgras{Eimer}
\textsuperscript{(9)} für das schmutzige Wasser.

%%K t06-c1-04-1 p
4. --- Im \VUgras{Vorraum} \textsuperscript{(10)} sehen wir ein paar
\VUgras{Kleiderhaken} \textsuperscript{(13)} und zwei \VUgras{Schirme}
\textsuperscript{(11)} im \VUgras{Schirmständer} \textsuperscript{(12)}.

%%K t06-c1-05-1 p
5. --- Auf der anderen Seite der Tür steht der \VUgras{Spiegelschrank}
\textsuperscript{(14)}, in dem die \VUgras{Wäsche} \textsuperscript{(15)} liegt.

%%K t06-c1-06-1 p
6. --- Während das \VUgras{Hausmädchen} \textsuperscript{(16)} das \VUgras{Bett}
\textsuperscript{(17)} macht, sehen wir uns seine Teile an. Das \VUgras{Bettgestell}
\textsuperscript{(20)} trägt eine gute \VUgras{Sprungfedermatratze}
\textsuperscript{(21)}, auf die Justine gleich eine wollene \VUgras{Matratze}
\textsuperscript{(22)} legen wird. Dann nimmt sie von den Stühlen die beiden
\VUgras{Laken} \textsuperscript{(23)} und die \VUgras{Decke} \textsuperscript{(24)}.
Dann legt sie ans Kopfende die \VUgras{Nackenrolle} \textsuperscript{(25)} und das
\VUgras{Kopfkissen} \textsuperscript{(26)}, die mit dem \VUgras{Federbett}
\textsuperscript{(27)} auf dem \VUgras{Bettvorleger} \textsuperscript{(32)} liegen.
Sie geht mit dem Staubwedel über die \VUgras{Vorhänge} \textsuperscript{(19)} und den
\VUgras{Betthimmel} \textsuperscript{(18)}, und ein Teil der Arbeit ist getan.

%%K t06-c1-07-1 p
7. --- Dann wird sie noch den \VUgras{Nachttisch} \textsuperscript{(28)} ein wenig
aufräumen, den \VUgras{Wecker} \textsuperscript{(29)} aufziehen, das
\VUgras{Nachtlicht} \textsuperscript{(30)} zurechtmachen und die \VUgras{Kerze}
\textsuperscript{(31)} wegnehmen, um den Leuchter zu putzen.

%%K t06-tit-2 sub
\VUsaut{5.0mm}
\VUpk{10.2pt}{40}{Der Nähkorb und das Kamingerät.}
\VUsaut{3.0mm}

%%K t06-c1-08-1 p
8. --- Auch der \VUgras{Nähkorb} \textsuperscript{(34)} neben dem \VUgras{Schemel}
\textsuperscript{(33)} muss dringend aufgeräumt werden. Sie wird die
\VUgras{Stickerei} \textsuperscript{(35)} zusammenlegen müssen, in den
\VUgras{Nähtisch} \textsuperscript{(38)} den \VUgras{Fingerhut}, das \VUgras{Garn}
\textsuperscript{(36)}, die \VUgras{Nadeln} \textsuperscript{(37)} und die
\VUgras{Schere} \textsuperscript{(39)} räumen und den Deckel des Tisches mit dem
\VUgras{Schlüssel} \textsuperscript{(40)} abschließen.

%%K t06-c1-09-1 p
9. --- Auf dem \VUgras{Rundtisch} \textsuperscript{(41)} in der Mitte wird Justine das
Wasser in der \VUgras{Karaffe} \textsuperscript{(42)} wechseln. Sie wird
\VUgras{Zucker} in die \VUgras{Zuckerdose} \textsuperscript{(43)} füllen, die
\VUgras{Zuckerzange} \textsuperscript{(44)} putzen und die \VUgras{Kleiderbürste}
\textsuperscript{(46)} wegräumen.

%%K t06-c1-10-1 p
10. --- Die rechte Seite des Zimmers macht ihr weniger Arbeit, weil die
\VUgras{Chaiselongue} \textsuperscript{(47)} und ihr \VUgras{Kissen}
\textsuperscript{(48)} am richtigen Platz stehen und weil bei dem \VUgras{Kamin}
\textsuperscript{(49)} nichts verrückt worden ist, da es nicht kalt ist.
\VUgras{Schaufel} \textsuperscript{(50)}, \VUgras{Zange} \textsuperscript{(51)},
\VUgras{Feuerböcke} \textsuperscript{(55)} und \VUgras{Aschenschirm}
\textsuperscript{(56)} sind sauber; der \VUgras{Ofenschirm} \textsuperscript{(54)}
steht unverrückt, und der \VUgras{Holzkasten} \textsuperscript{(52)} ist gut mit
\VUgras{Brennholz} \textsuperscript{(53)} gefüllt.

%%K t06-noto-1 noto
\VUnotes{143.2mm}{%
(*) Babo = ein kleines Kind; englisch \textit{baby}, französisch \textit{bébé},
italienisch \textit{bambino}.%
}

%%K t06-noto-2 noto
\VUnotes{143.2mm}{%
(*) Lambrequino = französisch \textit{lambrequin}, spanisch
\textit{lambrequines}, portugiesisch \textit{lambrequins}, und auch deutsch und
englisch \textit{lambrequins} --- also deutsch, englisch, französisch,
portugiesisch und spanisch gleichermaßen.%
}

%%K t06-c1-11-1 p
11. --- Sie wird noch die \VUgras{Kaminuhr} \textsuperscript{(57)}, die
\VUgras{Leuchter} \textsuperscript{(58)} und die \VUgras{Nippschalen}
\textsuperscript{(59)} abstauben und dann den \VUgras{Spiegel} \textsuperscript{(60)}
putzen.

%%K t06-c1-12-1 p
12. --- Die \VUgras{Wiege} \textsuperscript{(61)} des Babys (des babo (*)) ist schon
fertig.

%%K t06-tit-3 sub
\VUsaut{5.0mm}
\VUpk{10.2pt}{40}{Der Salon.}
\VUsaut{3.0mm}

%%K t06-c2-01-1 p
1. --- Jetzt der Salon. Es ist ein sehr schönes \VUgras{Zimmer}. Der Kamin in der
rechten Ecke ist mit einer zierlichen \VUgras{Statuette} \textsuperscript{(1)} und
zwei schönen \VUgras{Vasen} \textsuperscript{(3)} geschmückt und mit einer samtenen
\VUgras{Schabracke} \textsuperscript{(2)} (*) behängt.

%%K t06-c2-02-1 p
2. --- Damit die Funken vom Herd den Teppich nicht anzünden, ist ein kupferner
\VUgras{Funkenschirm} \textsuperscript{(4)} davorgestellt.

%%K t06-c2-03-1 p
3. --- Daneben steht offen ein eleganter \VUgras{Damensekretär}
\textsuperscript{(5)}. Dort schreibt die \VUgras{Hausherrin} \textsuperscript{(23)}
ihre Briefe.

%%K t06-c2-04-1 p
4. --- Das Fenster ist mit \VUgras{Stores} \textsuperscript{(6)} bekleidet, die von
\VUgras{Raffhaltern} \textsuperscript{(8)} an den \VUgras{Haken}
\textsuperscript{(7)} zurückgehalten werden, und mit schweren Stoffvorhängen, über
denen ein \VUgras{Volant} \textsuperscript{(9)} sitzt.

%%K t06-c2-05-1 p
5. --- In der Fensternische trägt ein \VUgras{Blumenständer} \textsuperscript{(11)}
eine grüne \VUgras{Pflanze} \textsuperscript{(10)}, eine prächtige Palme, deren
Blätter fast bis an die \VUgras{Petroleumlampe} \textsuperscript{(12)} reichen.

%%K t06-c2-06-1 p
6. --- Den \VUgras{Lampenschirm} \textsuperscript{(13)} hat Maria gemacht, die
reizende \VUgras{junge Dame} \textsuperscript{(14)} am \VUgras{Klavier}
\textsuperscript{(15)}. Sie sitzt sehr gerade auf dem \VUgras{Hocker}
\textsuperscript{(16)} und übt ein Stück. Sie ist sehr geschickt und sehr ordentlich,
wie ihr \VUgras{Notenschrank} \textsuperscript{(17)} zeigt: er ist in bester Ordnung.

%%K t06-tit-4 sub
\VUsaut{5.0mm}
\VUpk{10.2pt}{40}{Der Lausbub.}
\VUsaut{3.0mm}

%%K t06-c2-07-1 p
7. --- Ihr Bruder Paul sucht ein \VUgras{Buch} \textsuperscript{(19)} im
\VUgras{Bücherschrank} \textsuperscript{(18)}. Er ist ein \VUgras{Junge}
\textsuperscript{(20)}, der unruhig und ganz unerträglich ist. Er hängt sich an den
Schrank, um an das Buch zu kommen, das zu hoch steht, und läuft Gefahr, die
\VUgras{Büste} \textsuperscript{(21)} obendrauf herunterzureißen.

%%K t06-c2-08-1 p
8. --- Damit kommt er nur durch, weil seine Mutter gerade mit einer Freundin
spricht, einer \VUgras{Dame} \textsuperscript{(22)}, die für ein paar Tage zu Besuch
gekommen ist. Die Mutter sitzt auf dem \VUgras{Sofa} \textsuperscript{(24)} neben dem
\VUgras{Sessel} \textsuperscript{(25)}, und ihre Freundin auf dem \VUgras{Stuhl}
\textsuperscript{(27)}, von dem wir nur die \VUgras{Lehne} \textsuperscript{(26)}
sehen.

%%K t06-c2-09-1 p
9. --- Der große \VUgras{Rahmen} \textsuperscript{(30)} an der Wand enthält das
\VUgras{Bildnis} \textsuperscript{(29)} eines Verwandten. Im \VUgras{Album}
\textsuperscript{(32)} auf dem Tisch sind die \VUgras{Fotografien}
\textsuperscript{(33)} der übrigen Familie gesammelt. Die Kinder haben eben darin
geblättert: die \VUgras{Stühle} \textsuperscript{(31)} stehen noch um den Tisch
herum.

%%K t06-c2-10-1 p
10. --- Die porzellanene \VUgras{chinesische Vase} \textsuperscript{(34)} neben dem
\VUgras{Brieföffner} \textsuperscript{(35)} hat Maria zum Namenstag bekommen, und den
\VUgras{Wandschirm} \textsuperscript{(36)}, der die Ecke des Zimmers füllt, hat ihr
Onkel aus China mitgebracht.

%%K t06-c2-11-1 p
11. --- Mit den schönen \VUgras{Bildern} \textsuperscript{(37)} an der \VUgras{Wand}
\textsuperscript{(42)} und dem \VUgras{Kronleuchter} \textsuperscript{(38)} an der
Decke, dessen \VUgras{elektrische Lampen} \textsuperscript{(39)} abends so heiter
leuchten, sieht dieser Salon sehr freundlich aus. Der weiche \VUgras{Teppich}
\textsuperscript{(40)} auf dem Boden und die \VUgras{elektrische Klingel}
\textsuperscript{(41)}, die so praktisch ist, machen ihn vollkommen behaglich.

%%K t06-tit-5 sub
\VUsaut{3.6mm}
\VUpk{10.2pt}{40}{Das Esszimmer.}
\VUsaut{3.0mm}

%%K t06-c3-01-1 p
1. --- Es hat zum Essen geläutet. Die ganze Familie, außer Maria, sitzt um den
Tisch. Das Esszimmer wird von einem ausgezeichneten \VUgras{Kachelofen}
\textsuperscript{(1)} mit einem schlanken \VUgras{Ofenrohr} \textsuperscript{(2)}
angenehm gewärmt.

%%K t06-c3-02-1 p
2. --- In diesem Zimmer steht nicht viel. An der Wand, über dem \VUgras{Schemel}
\textsuperscript{(4)}, hängt ein zierlicher kleiner \VUgras{Wandbrunnen}
\textsuperscript{(3)}. Daneben steht die \VUgras{Anrichte} \textsuperscript{(5)}, auf
der die kalten Speisen, die Nachspeisen und das Geschirr abgestellt werden, das man
im Augenblick nicht braucht.

%%K t06-c3-03-1 p
3. --- Auf dem oberen Bord sehen wir daher ein \VUgras{Brathähnchen}
\textsuperscript{(7)}, einen \VUgras{Fisch} \textsuperscript{(8)} und eine
\VUgras{Schale} \textsuperscript{(10)} mit \VUgras{Obst} \textsuperscript{(9)}.
Darunter der \VUgras{Käse} \textsuperscript{(11)}, das \VUgras{Brot}
\textsuperscript{(12)} und die \VUgras{Menage} \textsuperscript{(13)}, die alles
enthält, was man für den Salat braucht: das Öl, den Essig, das \VUgras{Salz}
\textsuperscript{(14)} und den \VUgras{Pfeffer} \textsuperscript{(15)}. Auf dem
untersten Bord schließlich steht ein \VUgras{Kuchen} \textsuperscript{(17)} neben dem
\VUgras{Senftopf} \textsuperscript{(16)}.

%%K t06-c3-04-1 p
4. --- Die Uhr des Esszimmers, die man \VUgras{Wanduhr} \textsuperscript{(20)} nennt,
hat eine sehr ungewöhnliche Form. Sie sitzt in einem hölzernen Rahmen, der auch ein
\VUgras{Barometer} \textsuperscript{(19)} und ein \VUgras{Thermometer}
\textsuperscript{(18)} trägt.

%%K t06-tit-6 sub
\VUsaut{4.6mm}
\VUpk{10.2pt}{40}{Das Abendessen wird aufgetragen.}
\VUsaut{3.0mm}

%%K t06-c3-05-1 p
5. --- Rings um das Zimmer sind die Wände mit einer \VUgras{Holztäfelung}
\textsuperscript{(21)} aus gewachster Eiche verkleidet. Ein stoffener
\VUgras{Türvorhang} \textsuperscript{(22)} schließt die Tür zur Veranda leidlich ab.
Dorthin geht die Familie nach dem Essen zum Kaffee. Die \VUgras{Tassen}
\textsuperscript{(24)} und \VUgras{Untertassen} \textsuperscript{(25)} stehen schon
auf dem \VUgras{Geschirrschrank} \textsuperscript{(23)} bereit.

%%K t06-c3-06-1 p
6. --- Über dem Tisch ist eine \VUgras{Hängelampe} \textsuperscript{(26)} an der
Decke befestigt; ihr sehr helles Licht füllt abends das ganze Esszimmer.

%%K t06-c3-07-1 p
7. --- Sehen wir uns jetzt den \VUgras{Esstisch} \textsuperscript{(27)} selbst an. Als
die Familie hereinkam, war der Tisch gedeckt. Auf dem \VUgras{Tischtuch}
\textsuperscript{(28)} lagen an jedem Platz ein \VUgras{Teller}
\textsuperscript{(33)}, eine \VUgras{Serviette} \textsuperscript{(29)}, ein
\VUgras{Löffel} \textsuperscript{(34)}, eine \VUgras{Gabel} \textsuperscript{(35)},
ein \VUgras{Messer} \textsuperscript{(36)} und ein \VUgras{Glas}
\textsuperscript{(32)}. Dazu \VUgras{Flaschen} \textsuperscript{(30)} für den Wein
und eine \VUgras{Karaffe} \textsuperscript{(31)} für das Wasser.

%%K t06-c3-08-1 p
8. --- Der \VUgras{Diener} \textsuperscript{(39)} hat zuerst die
\VUgras{Suppenterrine} \textsuperscript{(37)} hereingebracht, in der noch Suppe
dampft. Jetzt trägt er den ersten \VUgras{Gang} \textsuperscript{(38)} auf, ein
Fleischgericht. Dann reicht er die schon genannten herum; und wenn die Herrschaft
nach dem \VUgras{Nachtisch} auf die Veranda hinübergegangen ist, räumt er den Tisch
ab und bringt das Esszimmer wieder in Ordnung.

%%K t06-c3-08-2 p
\textit{Anmerkung.} --- \textit{Die alltäglichen Wörter rund um das Essen sind hier
nur im Umriss angegeben. Die Liste wird auf den beiden Tafeln „Das Hotel“ (Nr. 13)
und „Der Markt“ (Nr. 15) vervollständigt.}

%%K t06-tit-7 sub
\VUsaut{4.6mm}
\VUpk{10.2pt}{40}{Die Küche.}
\VUsaut{3.0mm}

%%K t06-c4-01-1 p
1. --- In der Küche, die wir durch die halb offene Tür sehen, herrscht ein
malerisches Durcheinander. Vor dem \VUgras{Herdfeuer} \textsuperscript{(1)}, in dem
das \VUgras{Feuer} \textsuperscript{(3)} knistert, steht der \VUgras{Bratenwender}
\textsuperscript{(5)}. Ein schöner \VUgras{Braten} \textsuperscript{(7)} steckt am
\VUgras{Spieß} \textsuperscript{(6)} und wird gerade fertig.

%%K t06-c4-02-1 p
2. --- Der \VUgras{Blasebalg} \textsuperscript{(8)} und die \VUgras{Kohlenschaufel}
\textsuperscript{(9)}, die eben gebraucht wurden, liegen zusammen mit den
\VUgras{Scheiten} \textsuperscript{(10)} auf dem Boden.

%%K t06-c4-03-1 p
3. --- Der Kaminsims ist aufgeräumt: die \VUgras{Laterne} \textsuperscript{(11)}, der
\VUgras{Streichholzbehälter} \textsuperscript{(13)} mit den \VUgras{Streichhölzern}
\textsuperscript{(12)} darin, der \VUgras{Leuchter} \textsuperscript{(14)} und der
\VUgras{Handleuchter} \textsuperscript{(15)} mit seiner \VUgras{Kerze}
\textsuperscript{(16)} stehen alle an ihrem Platz. Aber der große
\VUgras{Kohlenkasten} \textsuperscript{(18)}, voll mit der \VUgras{Kohle}
\textsuperscript{(17)}, die die \VUgras{Köchin} \textsuperscript{(23)} eben aus dem
Keller geholt hat, ist mitten in der Küche stehen geblieben.

%%K t06-c4-04-1 p
4. --- Die arme Frau muss sich eben beeilen, wenn das Essen rechtzeitig fertig sein
soll. Gerade steht sie am \VUgras{Herd} \textsuperscript{(19)} und rührt eine
\VUgras{Soße} \textsuperscript{(24)} um, die nicht anbrennen darf.

%%K t06-c4-05-1 p
5. --- Im \VUgras{Backofen} \textsuperscript{(20)} bäckt ein Kuchen, den sie für die
Kinder gemacht hat. Über dem Herd hängen die \VUgras{Schöpfkelle}
\textsuperscript{(21)} und der \VUgras{Schaumlöffel} \textsuperscript{(22)}.

%%K t06-tit-8 sub
\VUsaut{4.0mm}
\VUpk{10.2pt}{40}{Das Küchengerät.}
\VUsaut{3.0mm}

%%K t06-c4-06-1 p
6. --- Die \VUgras{Reihe der Kasserollen} \textsuperscript{(25)}, die hinten im Raum
hängt, ist sehr sauber. Die schönen kupfernen \VUgras{Kasserollen}
\textsuperscript{(26)}, die \VUgras{Milchkanne} \textsuperscript{(27)}, das
\VUgras{Sieb} \textsuperscript{(28)} und die \VUgras{Reibe} \textsuperscript{(29)}
sind sorgfältig gescheuert.

%%K t06-c4-07-1 p
7. --- Das \VUgras{Salzfass} \textsuperscript{(30)} steht auf dem kleinen Bord neben
der \VUgras{Teekanne} \textsuperscript{(31)} und dem \VUgras{Ölballon}
\textsuperscript{(32)}. In der \VUgras{Schublade} \textsuperscript{(33)} liegen wohl
das Besteck und die Küchenmesser.

%%K t06-c4-08-1 p
8. --- Der \VUgras{Besen} \textsuperscript{(34)} und der \VUgras{Staubwedel}
\textsuperscript{(35)} stehen in einer Ecke. Die Köchin hat eben den
\VUgras{Hackklotz} \textsuperscript{(36)} gebraucht; sie hat ihn beim Fenster stehen
lassen.

%%K t06-c4-09-1 p
9. --- Ein \VUgras{Handtuch} \textsuperscript{(37)} hängt an der Wand neben dem
\VUgras{Trichter} \textsuperscript{(38)}. In diesem Teil der Küche werden der
\VUgras{Rost} \textsuperscript{(39)}, das \VUgras{Hackbeil} \textsuperscript{(40)}
und das \VUgras{Hackbrett} \textsuperscript{(41)} aufbewahrt. Über dem Hackbeil, auf
einem \VUgras{Bord} \textsuperscript{(42)}, stehen \VUgras{Töpfe}
\textsuperscript{(43)}, der \VUgras{Schmortopf} \textsuperscript{(44)} mit seinem
\VUgras{Deckel} \textsuperscript{(45)} und der \VUgras{Kessel}
\textsuperscript{(46)}. Die \VUgras{Bratpfanne} \textsuperscript{(47)} hängt gleich
daneben.

%%K t06-c4-10-1 p
10. --- Nur der kupferne \VUgras{Hahn} \textsuperscript{(48)} wirft einen Schimmer in
diese Ecke. Der \VUgras{Spülstein} \textsuperscript{(49)}, auf dem der große
\VUgras{Krug} \textsuperscript{(50)} steht, muss mit seinem kleinen Schrank sehr
praktisch sein: darin stehen das \VUgras{Essigfass} \textsuperscript{(51)} mit seinem
\VUgras{Hahn} \textsuperscript{(52)}, der \VUgras{Abfallkasten}
\textsuperscript{(53)} und das \VUgras{schmutzige Geschirr} \textsuperscript{(54)}.

%%K t06-tit-9 sub
\VUsaut{4.0mm}
\VUpk{10.2pt}{40}{Was sonst noch in der Küche steht.}
\VUsaut{3.0mm}

%%K t06-c4-11-1 p
11. --- Der große \VUgras{Zuber} \textsuperscript{(55)}, in dem das \VUgras{Gemüse}
\textsuperscript{(58)} gewaschen wird, und der gusseiserne \VUgras{Kessel}
\textsuperscript{(56)} auf dem \VUgras{Fliesenboden} \textsuperscript{(57)} gehören
wohl auch in diesen Schrank.

%%K t06-c4-12-1 p
12. --- An Feuerstellen fehlt es der Köchin gewiss nicht, denn hier auf der Ecke des
Tisches steht ein \VUgras{Gaskocher} \textsuperscript{(59)}, auf dem in einem kleinen
Topf Wasser heiß wird. Der \VUgras{Dampf} \textsuperscript{(60)}, der aufsteigt,
sagt uns, dass es kochen muss. Dieses Wasser ist wohl für den Kaffee, denn die
\VUgras{Kaffeekanne} \textsuperscript{(64)} und die \VUgras{Kaffeemühle}
\textsuperscript{(63)} stehen am anderen Ende des Tisches.

%%K t06-c4-13-1 p
13. --- Der Gaskocher ist an den \VUgras{Gasbrenner} \textsuperscript{(62)} mit
einem langen \VUgras{Gummischlauch} \textsuperscript{(61)} angeschlossen.

%%K t06-c4-14-1 p
14. --- Die \VUgras{Aufwartefrau} \textsuperscript{(66)}, die der Köchin zur
Hand geht, putzt das Gemüse für das Abendessen. Wenn es geschält und gewaschen ist,
legt sie es in die \VUgras{Kasserolle} \textsuperscript{(65)}, bis es Zeit zum Kochen
ist.

%%K t06-tit-10 sub
\VUsaut{4.0mm}
{\centering\fontsize{10.2pt}{10.2pt}\selectfont\textls[40]{\scshape Das Badezimmer.} \textit{(Der Baderaum.)}\par}
\VUsaut{3.0mm}

%%K t06-c5-01-1 p
1. --- Noch ein letztes Mal neugierig sein, und wir kennen die wichtigsten Zimmer
des Hauses: die Einrichtung des Badezimmers. Das dauert nicht lange, denn der Raum
ist klein, und wir sehen auf einen Blick, dass die \VUgras{Badewanne}
\textsuperscript{(1)} einen guten \VUgras{Badeofen} \textsuperscript{(2)} hat, dass
die \VUgras{Seifenschale} \textsuperscript{(3)} in Reichweite des Badenden liegt und
dass der \VUgras{Handtuchhalter} \textsuperscript{(5)} das Trocknen der
\VUgras{Handtücher} \textsuperscript{(4)} leicht machen muss.

%%K t06-c5-02-1 p
2. --- Die Wände dieses Raums sind mit \VUgras{glasierten Kacheln}
\textsuperscript{(6)} verkleidet, so dass man eine \VUgras{Dusche}
\textsuperscript{(7)} einbauen konnte, ohne dass das umherspritzende Wasser ihnen
schadet. Ein kleines zinkernes \VUgras{Becken} \textsuperscript{(8)} für Fußbäder
vervollständigt die Einrichtung.
\VUsaut{6.0mm}
\VUfilet{22mm}
